\section{Reduction by the number of reflections}\label{sec:reflection-count}

Let \(S\) be an arbitrary sequence over \(G=\Dih(A)\) of length
\begin{equation}\label{eq:ambient-length}
  |S|=2Q+D_A.
\end{equation}
Let \(a\) be the number of reflections in \(S\), and put
\[
  b=2Q+D_A-a
\]
for the number of rotations. For a subgroup \(K\le A\) and a sequence
\(Y\) over \(G\), let \(m_K(Y)\) denote the number of rotations
\((x,0)\) of \(Y\) with \(x\in K\).

\subsection{At most one reflection}

\begin{proposition}\label{prop:few-reflections}
If \(a\le1\), then \(S\) contains a product-one subsequence of length
\(2Q\) consisting entirely of rotations.
\end{proposition}

\begin{proof}
There are at least
\[
  b\ge 2Q+D_A-1=2Q+\D(A)-1
\]
rotations. Apply \Cref{thm:GMO}, in its existence form
\eqref{eq:GMO-existence}, to the rotation sequence over \(A\), with
weight set \(\{1\}\) and target \(n=2Q\). The target satisfies
\(2Q\ge |A|\), and the displayed inequality is the required length
hypothesis. Since \(\exp(A)\mid Q\), one has \(2QA=\{0\}\). Thus
there are exactly \(2Q\) rotations with ordinary sum zero.
\end{proof}

\subsection{Many reflections}

\begin{proposition}\label{prop:many-reflections}
If \(a\ge D_A+1\), then \(S\) contains a product-one subsequence of
length \(2Q\).
\end{proposition}

\begin{proof}
Pair reflections with reflections and rotations with rotations. Since
\(2Q+D_A\) is odd, exactly one term is left unpaired, and the number of
pairs is
\[
  \left\lfloor\frac a2\right\rfloor
  +\left\lfloor\frac b2\right\rfloor
  =Q+\frac{D_A-1}{2}.
\]
Assign to a reflection pair \((x,1)\cdot(y,1)\) the label \(x-y\), and
to a rotation pair \((u,0)\cdot(v,0)\) the label \(u+v\). By
\Cref{prop:plus-minus-bound}, the sequence of pair labels has length at
least
\[
  Q+\Dpm(A)-1.
\]
Apply \eqref{eq:GMO-existence} in \(A\), with weight set
\(\{-1,1\}\) and target \(Q=|A|\). Since \(QA=\{0\}\), we obtain
exactly \(Q\) pairs and a sign for each selected pair such that the
signed sum of their labels is zero.

At most \(\lfloor b/2\rfloor\le Q-1\) of all pairs are rotation pairs,
because \(a\ge D_A+1\). Hence at least one selected pair is a
reflection pair. For a reflection pair with selected weight
\(\varepsilon\), assign signs \(\varepsilon\) and \(-\varepsilon\)
to its two reflections. For a rotation pair, assign the selected weight
to both rotations. The selected reflection signs are balanced, and the
total signed vector sum is zero. By \Cref{lem:sign-realization}, the
\(2Q\) selected terms have a product-one ordering.
\end{proof}

\subsection{The intermediate range}

Assume now that
\begin{equation}\label{eq:middle-range}
  2\le a\le D_A.
\end{equation}
Let \(e\) be the largest even integer not exceeding \(a\), and put
\begin{equation}\label{eq:middle-target}
  \ell=2Q-e.
\end{equation}
Then \(e>0\), \(\ell\ge Q\), and \(\ell\le b\). Moreover,
\begin{equation}\label{eq:middle-surplus}
  b-\ell=
  \begin{cases}
    D_A,& a\text{ even},\\
    D_A-1,& a\text{ odd},
  \end{cases}
  \qquad
  b-\ell\ge \Dpm(A)-1.
\end{equation}
Indeed, \(e\le D_A\le Q\), and the final inequality follows from
\Cref{prop:plus-minus-bound}.

\begin{proposition}[Intermediate-range reduction]
\label{prop:middle-reduction}
Under \eqref{eq:middle-range}, either \(S\) contains a product-one
subsequence of length \(2Q\), or there is a nonzero proper subgroup
\(K<A\) such that
\begin{equation}\label{eq:concentration-K}
  m_K(S)\ge b-[A:K]+2.
\end{equation}
\end{proposition}

\begin{proof}
Apply the structural part of \Cref{thm:GMO} to the sequence of all
\(b\) rotation vectors, with ambient group \(A\), weight set
\(\{-1,1\}\), and target \(\ell\). The target condition
\(\ell\ge|A|\) and the length surplus are exactly
\eqref{eq:middle-surplus}.

If \(\Sigma_\ell^{\{-1,1\}}(S_{\rm rot})=A\), select any \(e\)
reflections and assign half of them each sign. Let their signed vector
sum be \(t\in A\). Fullness supplies \(\ell\) rotations with a signed
sum of \(-t\). The two selections are disjoint, and their total length
is \(e+\ell=2Q\). The sign assignment on the reflections is balanced,
so \Cref{lem:sign-realization} gives a product-one ordering.

Otherwise, \Cref{thm:GMO} gives a proper subgroup \(K<A\) and a
rotation subsequence of length at least \(b-[A:K]+2\) satisfying
\eqref{eq:GMO-weight-coset}. For every one of its values \(x\), the
elements \(x\) and \(-x\) lie in a common \(K\)-coset. Hence
\(2x\in K\). Since \(A/K\) has odd order, multiplication by two is
injective on \(A/K\), and therefore \(x\in K\). This proves
\eqref{eq:concentration-K}.

It remains to exclude \(K=0\) from the second alternative. In that
case at least
\[
  b-Q+2=Q+D_A-a+2\ge D_A
\]
rotations are the identity, so \Cref{lem:identity-padding} gives the
required subsequence directly.
\end{proof}
